\section{Introduction}

Educational settings today face ongoing challenges related to student motivation, engagement,
and effective learning.
Factors such as increased academic pressure, digitalization, and the
growing presence of AI-based tools have contributed to a more complex learning environment
for both students and teachers.

This project explores how learning outcomes in education can be enhanced through the design
and development of a software system aimed at improving student engagement, participation,
and motivation.

The work is grounded in user-centered design, with insights derived from
interviews, personas, and domain modeling.

This report combines and extends work from previous assignments, documenting both the
pre-study of the problem domain and the iterative development of a potential solution.
It includes analysis of user needs, domain modeling, requirements, implementation choices,
and reflections on development practices such as Scrum, Extreme Programming (XP), and
continuous integration.

The project follows principles from design science and design thinking, emphasizing the
importance of understanding the problem before proposing solutions, and iteratively refining
both the problem and the system based on feedback and reflection.
The purpose of this report is therefore to present a structured overview of the project’s progression,
key decisions, and lessons learned throughout the development process.
